\section{Block-Isoclinic Form}
\label{sec:isoclinic}

The inquiry into simultaneous arithmetic across several cells has a purely
linear-algebraic consequence.  Cells belonging to one factor node share the
same angle, so a block is a repeated rotation on mutually orthogonal
two-planes.

Historically, the June 10--11 block/SIMD inquiry produced exactly the invariant
proved below: equal-angle lanes carry an $O(w)$-equivariant representation of
one isoclinic rotation.  The statement is independent of any instruction set.

Let $R_\theta=\left[\begin{smallmatrix}c&-s\\s&c\end{smallmatrix}\right]$ and
let $w$ be the number of pairs in a node.  After a fixed coordinate shuffle,
the rotation portion is
\begin{equation}
 \mathcal{R}_{\theta,w}=I_w\otimes R_\theta.
 \label{eq:isoclinic-block}
\end{equation}
The add-subtract portion is a signed direct sum of $H_2$ maps.

\begin{proposition}[Lane equivariance]
For every $U\in O(w)$,
\begin{equation}
 (U\otimes I_2)\mathcal{R}_{\theta,w}
 =\mathcal{R}_{\theta,w}(U\otimes I_2).
 \label{eq:lane-equivariance}
\end{equation}
Moreover, the complete block is $\sqrt2$ times an orthogonal map.
\end{proposition}
\begin{proof}
The mixed-product rule for Kronecker products gives both sides as
$U\otimes R_\theta$.  The Gram identity follows from
$R_\theta^{\trans}R_\theta=I_2$ and the Hadamard identity.
\end{proof}

Thus a collection of lanes is a representation of one algebraic operator, not
a different Fourier transform.  The important invariant is the common
two-plane angle.  DIF associates that angle with a factor-tree node, DIT with a
spectral position, and DIP with a diagonal phase span.
